\documentclass[11pt,a4paper]{article}

\usepackage[T1]{fontenc}
\usepackage[utf8]{inputenc}
\usepackage{lmodern}
\usepackage[a4paper,top=2.4cm,bottom=2.2cm,left=2.2cm,right=2.2cm]{geometry}
\usepackage{amsmath}
\usepackage{amssymb}
\usepackage{graphicx}
\graphicspath{{./}}
\usepackage{array}
\usepackage{makecell}
\renewcommand{\arraystretch}{1.25}
\usepackage{enumitem}
\usepackage{fancyhdr}
\usepackage{xcolor}

\newcommand{\ohms}{\ensuremath{\Omega}}
\newcommand{\degC}{\ensuremath{^{\circ}}C}
\newcommand{\dg}{\ensuremath{^{\circ}}}
\newcommand{\pow}[1]{\ensuremath{^{#1}}}

\newcommand{\topiclabel}[1]{{\hfill\normalfont\small\itshape\textcolor{black!55}{#1}}}

% Per-question head: \examq{paper-id}{number}{topics}{marks}
% Renders: paper-id - Q<number> - topics - marks  (one line, same font).
% Same command and same rendered line as the answer paper.
\newcommand{\swmarks}{}
\newcommand{\examq}[4]{%
  \gdef\swmarks{#4}%
  \par\addvspace{16pt}\noindent
  {\large\bfseries #1 - Q#2 - #3 - #4}\par\vspace{5pt}%
}
% \total now takes NO argument: it prints the marks stored by \examq.
\newcommand{\total}{\par\vspace{2mm}\noindent\hspace*{\fill}[Total:\ \swmarks]\par\vspace{2mm}}

\newlength{\anslineskip}
\setlength{\anslineskip}{8mm}

\newcounter{ansln}
\newcommand{\Alines}[2]{%
  \par\vspace{1mm}%
  \setcounter{ansln}{1}%
  \loop\ifnum\value{ansln}<#1\relax
    \noindent\mbox{}\dotfill\par\vspace{\anslineskip}%
    \stepcounter{ansln}%
  \repeat
  \noindent\mbox{}\dotfill\hspace{0.6em}\mbox{#2}\par
  \vspace{1mm}%
}

\newcommand{\plainline}{\par\noindent\mbox{}\dotfill\par\vspace{\anslineskip}}

\newcommand{\labelline}[1]{\par\vspace{1mm}\noindent #1\hspace{0.5em}\dotfill\par\vspace{\anslineskip}}

\newcommand{\ansval}[3]{%
  \par\vspace{5mm}%
  \noindent\hspace*{3.5cm}#1\hspace{0.5em}\dotfill\hspace{0.5em}#2\hspace{0.5em}\makebox[1.1cm][r]{#3}\par
  \vspace{3mm}%
}

\renewcommand{\markright}[1]{\par\nopagebreak\noindent\hspace*{\fill}\mbox{#1}\par\vspace{1mm}}

\newcommand{\qfig}[3][0.5\linewidth]{%
  \par\vspace{2mm}
  \begin{center}%
    \IfFileExists{#2}%
      {\includegraphics[width=#1]{#2}}%
      {\fbox{\parbox[c][26mm][c]{#1}{\centering\ttfamily\footnotesize #2}}}\\[4pt]
    \textbf{#3}%
  \end{center}%
  \vspace{2mm}\par
}

\newlist{parts}{enumerate}{1}
\setlist[parts]{label=\textbf{(\alph*)},leftmargin=2.1em,labelsep=0.6em,itemsep=10pt,topsep=6pt,parsep=4pt}

\newlist{subparts}{enumerate}{1}
\setlist[subparts]{label=\textbf{(\roman*)},leftmargin=2.3em,labelsep=0.6em,itemsep=10pt,topsep=6pt,parsep=4pt}

\pagestyle{fancy}
\fancyhf{}
\fancyhead[L]{\footnotesize 4037/11/M/J/25}
\fancyhead[C]{\footnotesize Cambridge O Level Additional Mathematics -- Paper 1}
\fancyhead[R]{\footnotesize May/June 2025}
\fancyfoot[C]{\footnotesize\thepage}
\renewcommand{\headrulewidth}{0.4pt}
\renewcommand{\footrulewidth}{0pt}

\setlength{\parindent}{0pt}

\begin{document}



%==================== QUESTION 1 ====================
\examq{4037/11/M/J/25}{1}{Vectors in Two Dimensions}{5}
\begin{parts}
  \item Given that $\overrightarrow{PQ}=\begin{pmatrix}-3\\7\end{pmatrix}$ and $4\overrightarrow{PR}=\begin{pmatrix}-2\\8\end{pmatrix}$, find $\overrightarrow{RQ}$.
  \markright{[2]}
  \vspace{5.5cm}

  \item The vectors $\mathbf{a}$, $\mathbf{b}$ and $\mathbf{c}$ are such that $\mathbf{a}=\alpha\mathbf{i}+6\mathbf{j}$, $\mathbf{b}=4\mathbf{i}+\beta\mathbf{j}$ and $\mathbf{c}=(2\alpha+5\beta)\mathbf{i}+20\mathbf{j}$, where $\alpha$ and $\beta$ are scalars.

  Given that $\mathbf{c}=3\mathbf{a}-2\mathbf{b}$, find the values of $\alpha$ and $\beta$.
  \markright{[3]}
  \vspace{10cm}
\end{parts}
\total

%==================== QUESTION 2 ====================
\examq{4037/11/M/J/25}{2}{Quadratic Functions}{4}
Solve the inequality $(3-x)(5x+8)\geqslant 9-3x$.
\markright{[4]}
\vspace{18cm}
\total

%==================== QUESTION 3 ====================
\examq{4037/11/M/J/25}{3}{Coordinate Geometry of the Circle}{6}
Point $A$ has coordinates $(3,-1)$.

A circle has equation $\left(x-4\right)^{2}+\left(y+3\right)^{2}=5$.
\begin{parts}
  \item Show that $A$ lies on the circumference of the circle.
  \markright{[1]}
  \vspace{2.5cm}

  \item Given that $AB$ is a diameter of the circle, find the coordinates of $B$.
  \markright{[2]}
  \vspace{8.5cm}

  \item Find the equation of the tangent to the circle at $A$.
  \markright{[3]}
  \vspace{8cm}
\end{parts}
\total

%==================== QUESTION 4 ====================
\examq{4037/11/M/J/25}{4}{Quadratic Functions \textperiodcentered\ Simultaneous Equations \textperiodcentered\ Logarithmic \& Exponential Functions}{9}
\begin{parts}
  \item Solve the equation $x^{\frac{1}{3}}-x^{\frac{1}{6}}=2$.
  \markright{[4]}
  \vspace{7cm}

  \item Solve the simultaneous equations
  \[
  \begin{aligned}
    \lg(x+2y)&=0\\
    x^{2}+4xy+y&=1\,.
  \end{aligned}
  \]
  \markright{[5]}
  \vspace{10cm}
\end{parts}
\total

%==================== QUESTION 5 ====================
\examq{4037/11/M/J/25}{5}{Equations, Inequalities \& Graphs \textperiodcentered\ Trigonometry}{5}
\begin{parts}
  \item Solve the equation $5\left|2x-1\right|+8=23$.
  \markright{[3]}
  \vspace{4.5cm}

  \item On the axes, sketch the graph of $y=5\sin x-2$ for $0\dg\leqslant x\leqslant 360\dg$.
  \markright{[2]}
  \qfig[0.72\linewidth]{4037_11_M_J_25_fig1.png}{}
\end{parts}
\total

%==================== QUESTION 6 ====================
\examq{4037/11/M/J/25}{6}{Differentiation}{8}
A curve has equation $y=\left(\dfrac{x^{2}-1}{x^{2}+1}\right)^{4}$.
\begin{parts}
  \item Show that $\dfrac{\mathrm{d}y}{\mathrm{d}x}$ can be written as $\dfrac{Ax\left(x^{2}-1\right)^{3}}{\left(x^{2}+1\right)^{5}}$, where $A$ is a positive integer to be found.
  \markright{[5]}
  \vspace{17cm}

  \item \begin{subparts}
    \item Show that the curve has stationary points where $x=-1$, $x=0$ and $x=1$.
    \markright{[1]}
    \vspace{3cm}

    \item Use the first derivative test to determine which two stationary points have the same nature and state whether they are maximum or minimum points.
    \markright{[2]}
    \vspace{14cm}
  \end{subparts}
\end{parts}
\total

%==================== QUESTION 7 ====================
\examq{4037/11/M/J/25}{7}{Factors of Polynomials}{6}
\textbf{Solutions to this question by accurate drawing will not be accepted.}

Find the $x$-coordinates of the points where the curve $y=(2x-9)\left(x^{2}+5\right)+42$ cuts the $x$-axis.
\markright{[6]}
\vspace{18cm}
\total

%==================== QUESTION 8 ====================
\examq{4037/11/M/J/25}{8}{Logarithmic \& Exponential Functions}{7}
\begin{parts}
  \item Write down the set of values of $x$ for which $\log_{5}(12x-4)$ exists.
  \markright{[1]}
  \vspace{2.5cm}

  \item Solve the equation $\log_{5}(12x-4)=\dfrac{6}{\log_{x}125}+1$.
  \markright{[6]}
  \vspace{16cm}
\end{parts}
\total

%==================== QUESTION 9 ====================
\examq{4037/11/M/J/25}{9}{Differentiation \textperiodcentered\ Integration}{10}
\qfig[0.6\linewidth]{4037_11_M_J_25_fig2.png}{}
The point $A$ with $x$-coordinate 2 lies on the curve $y=\sqrt{4x+1}$.\\
The diagram shows part of this curve and the tangent to the curve at $A$.

Find the area of the shaded region enclosed by the curve, the tangent and the $x$-axis.
\markright{[10]}
\vspace{8cm}

\newpage
\textbf{Continuation of working space for Question 9.}
\vspace{17cm}
\total

%==================== QUESTION 10 ====================
\examq{4037/11/M/J/25}{10}{Trigonometry}{7}
\begin{parts}
  \item Given that $0\leqslant\theta<\dfrac{\pi}{2}$, show that $\dfrac{\sin\theta}{\sqrt{\operatorname{cosec}^{2}\theta-1}}+\dfrac{1}{\sqrt{1+\tan^{2}\theta}}$ can be written as $\sec\theta$.
  \markright{[4]}
  \vspace{9cm}

  \item Given that $\sec x=\alpha$, where $\dfrac{3\pi}{2}<x\leqslant 2\pi$, find $\sin x$ in terms of $\alpha$.
  \markright{[3]}
  \vspace{9cm}
\end{parts}
\total

%==================== QUESTION 11 ====================
\examq{4037/11/M/J/25}{11}{Arithmetic \& Geometric Progressions}{9}
An arithmetic progression has common difference $d$.\\
The 3rd term of this progression is 10.
\begin{parts}
  \item Write down expressions for the 1st term and the 2nd term of this progression.\\
  Give your answers in terms of $d$ only.
  \markright{[2]}
  \vspace{5cm}

  \item When each of the first 3 terms is squared, the sum of these squares is 140 .\\
  There are two possible values for $d$.

  Using your answer to part \textbf{(a)}, find the sum of the first 200 terms of the progression with the smaller value of $d$.
  \markright{[7]}
  \vspace{12cm}
\end{parts}
\total

\begin{center}\textbf{Question 12 is printed on the next page.}\end{center}

%==================== QUESTION 12 ====================
\examq{4037/11/M/J/25}{12}{Permutations \& Combinations}{4}
In this question $n\geqslant 6$ .\\
Use an algebraic method to show that $^{n}\mathrm{C}_{5}-{}^{n-1}\mathrm{C}_{5}$ can be written as $^{n-1}\mathrm{C}_{4}$.
\markright{[4]}
\vspace{12cm}
\total

\end{document}